\scanpage{76}
% archive sheet number: 75
% L1
\[
\psi^{(u)}\psi^{(u^2)}\cdots\psi^{(u^{n-1})}=1
\]
% L2-5
(où $\psi^{(u^k)}=u^k\psi u^{k-1}$), d'où on conclut, en désignant par $\widetilde\psi$ le prolongement de $\psi$ à $G^\Omega$, et notant que les $\psi^{(u^k)}$ sont encore des $\unsure{k|L}$-automorphismes \inserted{du gr. algébrique} $G^L$ par la \deleted{\unsure{restriction}}\inserted{restriction} des automorphismes $\widetilde\psi^{(\widetilde u^k)}$ de $G^\Omega$ :
% L6
\[
\widetilde\psi^{\widetilde u}\widetilde\psi^{\widetilde u^2}\cdots\widetilde\psi^{\widetilde u^{n-1}}=1 \tag*{(*)}
\]
% L7-8
Soit $\unsure{\varphi|\phi}\in Z^1(\Gamma,G^{L\psi})$, donné par $\varphi(u)=a\in G^{L\unsure{\psi}}$, avec
\[
a\,a^{\unsure{\sigma|v}}a^{\sigma^2}\cdots a^{\sigma^{n-1}}=1 \tag*{(**)}
\]
% L9-11
On cherche $b\in G^L$ tel que $a=b(b^\sigma)^{-1}$ ; on va voir qu'il suffit de \unsure{y|à} trouver $b\in G^\Omega$ tel que $a=b(b^{\widetilde\sigma})^{-1}$ : on va voir qu'on aura automatiquement $b\in G^L$.
% L12-15
Pour ceci, on \deleted{\unsure{note que}} \deleted{\unsure{part de}\ $a=b(b^{\widetilde\sigma})^{-1}$} : \unsure{dans} $(**)$
\[
\deleted{\unsure{a^{\sigma^k}|a^{u^k}}=a^{\widetilde\sigma^k}}
\]
\[
\deleted{\widetilde\sigma^k=(\widetilde\psi\widetilde u)^k=\widetilde\psi\widetilde\psi^{\widetilde u}\widetilde\psi^{\widetilde u^2}\cdots\widetilde\psi^{\widetilde u^{\unsure{k|k-1}}}\widetilde u^k}
\]
\deleted{(récurrence sur $k$)}
\deleted{D'où}
\[
\deleted{a^{\sigma^k}=a^{\widetilde\sigma^k}=(\widetilde\psi\cdots\widetilde\psi^{\widetilde u^{k-1}})\,b\,b^\sigma}
\]
\[
\deleted{\unsure{b^{\widetilde\sigma-1}|b^{\widetilde\sigma^{-1}}}}.
\]
% L16-18
en notant que
\[
a^{\sigma^k}=a^{\widetilde\sigma^k}=b^{\widetilde\sigma^k}(b^{\widetilde\sigma^{k+1}})^{-1}.
\]
d'où on tire
\[
b(b^{\widetilde\sigma^n})^{-1}=1\qquad\text{soit}\qquad b^{\widetilde\sigma^n}=b
\]
% L19-21
Or on a
\[
\widetilde\sigma^k=(\widetilde\psi\widetilde u)^k=\widetilde\psi\widetilde\psi^{\widetilde u}\cdots\widetilde\psi^{\widetilde u^{k-1}}\widetilde u^k
\]
(récurrence sur $k$) d'où
\[
\widetilde\sigma^n=(\widetilde\psi\widetilde\psi^{\widetilde u}\cdots\widetilde\psi^{\widetilde u^{n-1}})\widetilde u^n=\widetilde u^n
\]
% L22-24
en vertu de $(*)$, on a donc $b^{\unsure{\widetilde u^n}}=b$, i.e. $b\in G^L$.

On est maintenant ramené au théorème géométrique suivant :
